\subsubsection{Empreinte mémoire selon la longueur et la divergence}
À longueur fixée, WFA2 explore davantage de diagonales lorsque la distance augmente. Dans la version dense initiale de NvPD, la taille des tampons dépend de la longueur des séquences. Le calcul du score place la plus longue chaîne en colonne. Sur l'alphabet à quatre bases, le GPU réserve quatre lignes d'indices et deux lignes de scores, chacune en entiers de quatre octets. Pour des séquences de longueurs \(m,n\), cette allocation vaut
\[
  M_{\mathrm{device}}=4(4+2)\bigl(\max(m,n)+1\bigr)\quad [\mathrm{bytes}].
\]
Le panneau a) de la figure~\ref{fig:synthetic-memory} porte sur les quinze paires de chaque longueur. Le RSS est le pic résident du processus hôte. Les panneaux b) à e) utilisent la première paire de chaque taux. Pour les deux bibliothèques sur Grace, un second processus relève le RSS après lecture et pendant les appels, avec un intervalle cible d'une milliseconde. Le maximum ainsi échantillonné est une borne inférieure du pic de mémoire.

\chart{comparison/memoire_synthetique.pdf}{Le panneau a) compare les tampons CUDA au RSS hôte de NvPD (échelle logarithmique). Il donne la moyenne et les extrêmes sur quinze paires distinctes par longueur. Les panneaux b) à e) donnent les moyennes des maxima échantillonnés pour Edlib et WFA2. Les traits épais indiquent un écart-type, les traits fins l'IC à 95\,\% sur neuf processus à 5 et 20~knt et cinq aux deux autres longueurs. Les échelles horizontales sont propres à chaque longueur.}{fig:synthetic-memory}{python/studies/synthetic_memory.py}
À 1~Mnt, les tampons du noyau initial occupent de 24\,000\,024 à 24\,016\,800 octets sur les quinze paires. Les insertions et les délétions changent leur taille lorsque la chaîne mutée dépasse la référence. Le pic RSS hôte de NvPD vaut 216\,Mio pour chacune de ces paires, contre 180\,Mio de 5 à 100~knt.

Avec le calcul en bande, un groupe de 32 colonnes occupe 32 octets dans l’index. Les quatre masques de présence et les quatre positions précédant le groupe permettent de retrouver la dernière occurrence de chaque base. Les scores occupent deux tranches dont la largeur dépend de la borne \(U\). Le panneau a) de la figure~\ref{fig:nvpd-focus-memoire} compare les allocations explicites des deux noyaux.

Le panneau b) mesure l’effet du nombre de processus. Sur un même GH200, cinq paires de 100\,000 bases à 10\,\% sont traitées successivement par un processus, puis comparées à une exécution concurrente avec un processus par paire. L’ordre des deux modes varie entre les cinq blocs. Après échauffement, Linux attribue à chaque processus une part \(1/k\) de chaque page partagée par \(k\) processus. La somme de ces mémoires résidentes proportionnelles (PSS) décrit la place du lot dans la mémoire hôte.
\chart{comparison/nvpd_focus_memory.png}{À gauche, les tampons GPU sont mesurés sur cinq paires distinctes par longueur à 2\,\%. À droite, la PSS hôte du lot est relevée après échauffement. Les durées sous les libellés donnent la moyenne et l’écart-type sur cinq blocs. Les traits fins indiquent les intervalles à 95\,\%.}{fig:nvpd-focus-memoire}{bench/repeated/divergence/focus_analysis.py}

À un million de bases et 2\,\%, le pic RSS hôte atteint en moyenne 159,5\,Mio avec NvPD en bande, 10,1\,Mio avec WFA2 et 9,6\,Mio avec Edlib. Le lot à processus unique prend 0,365\,s, contre 0,389\,s avec cinq processus.

\synthese{À un million de bases et 2\,\%, la bande ramène les tampons GPU de 22,89 à 1,59\,Mio par paire. Pour le lot de cinq paires de 100\,000 bases, un seul processus ramène la PSS hôte de 726,2 à 154,9\,Mio.}
